
\documentclass{article}
\usepackage{colortbl}
\usepackage{makecell}
\usepackage{multirow}
\usepackage{supertabular}

\begin{document}

\newcounter{utterance}

\centering \large Interaction Transcript for game `codenames', experiment `risk\_low', episode 3 with deepseek{-}v4{-}pro.
\vspace{24pt}

{ \footnotesize  \setcounter{utterance}{1}
\setlength{\tabcolsep}{0pt}
\begin{supertabular}{c@{$\;$}|p{.15\linewidth}@{}p{.15\linewidth}p{.15\linewidth}p{.15\linewidth}p{.15\linewidth}p{.15\linewidth}}
    \# & $\;$A & \multicolumn{4}{c}{Game Master} & $\;\:$B\\
    \hline

    \theutterance \stepcounter{utterance}  
    & & & \multicolumn{2}{p{0.3\linewidth}}{
        \cellcolor[rgb]{0.9,0.9,0.9}{
            \makecell[{{p{\linewidth}}}]{
                \texttt{\tiny{[GM$|$GM]}}
                {'hidden': {'team': ['χάρη', 'μάτι', 'χιονάνθρωπος', 'σύνδεσμος', 'συμβολοσειρά', 'γκρεμός', 'κάλτσα', 'παγωτό', 'λιμάνι'], 'innocent': ['αεροπλάνο', 'λιμουζίνα', 'προσωπικό', 'μπότα', 'κατηγορώ', 'καρέκλα', 'μαστίγιο', 'γεράκι'], 'opponent': ['χαλκός', 'κύτταρο', 'σχολείο', 'πλυντήριο', 'ζώνη', 'εργαστήριο', 'γραμμή', 'αστυνομία'], 'assassin': []}, 'revealed': {'team': {'team': [], 'innocent': [], 'opponent': [], 'assassin': []}, 'opponent': {'team': [], 'innocent': [], 'opponent': [], 'assassin': []}}}
            }
        }
    }
    & & \\ \\

    \theutterance \stepcounter{utterance}  
    & & \multicolumn{4}{p{0.6\linewidth}}{
        \cellcolor[rgb]{0.9,0.9,0.9}{
            \makecell[{{p{\linewidth}}}]{
                \texttt{\tiny{[P1$\langle$GM]}}
                \texttt{Δώστε μία μονολεκτική ένδειξη που σχετίζεται με κάποιες από τις λέξεις της ομάδας σας από τον παρακάτω κατάλογο. Μπορείτε να επιλέξετε σε πόσες και σε ποιες λέξεις αναφέρεται η ένδειξή σας.} \\
\texttt{Ο στόχος είναι να μαντέψει ο συνεργάτης σας όλες τις λέξεις της ομάδας σας με βάση τις ενδείξεις που του δίνετε, πριν η άλλη ομάδα αποκαλύψει όλες τις δικές της. Αν ο συνεργάτης σας μαντέψει μια λέξη του 'αντιπάλου' ή μια λέξη 'αντιπερισπασμού', η σειρά μαντέματός του τελειώνει και η αντίπαλη ομάδα παίρνει τη σειρά της. Αν ο συνεργάτης σας μαντέψει μια λέξη 'θανάτου', το παιχνίδι τελειώνει και η ομάδα σας χάνει αμέσως.} \\
\texttt{Όσο περισσότερες λέξεις σχετίζονται με την ένδειξή σας, τόσο πιο γρήγορα μπορείτε να κερδίσετε το παιχνίδι, αλλά τόσο πιο δύσκολοι μπορεί να είναι οι συσχετισμοί για να τους μαντέψει ο συνεργάτης σας, οπότε επιλέξτε τη λέξη{-}ένδειξή σας με σύνεση. Η λέξη{-}ένδειξη πρέπει να σχετίζεται σημασιολογικά με τις λέξεις{-}στόχους, δεν μπορεί να είναι μία από τις λέξεις των καταλόγων ή να περιέχει μέρη τους.} \\
\\ 
\texttt{Να δίνετε πάντα τη μονολεκτική σας ένδειξη και τον κατάλογο των σχετικών λέξεων{-}στόχων χωρισμένο με κόμματα με την ακόλουθη μορφή και να κάνετε τις απαντήσεις σας όσο το δυνατόν πιο σύντομες, μην περιλαμβάνετε ποτέ άλλο κείμενο από αυτό που απαιτείται σε αυτή τη μορφή:} \\
\\ 
\texttt{ΥΠΟΔΕΙΞΗ: <ΛΕΞΗ>} \\
\texttt{ΣΤΟΧΟΙ: <ΣΤΟΧΟΙ>} \\
\\ 
\texttt{Οι λέξεις της ομάδας σας είναι: χάρη, μάτι, χιονάνθρωπος, σύνδεσμος, συμβολοσειρά, γκρεμός, κάλτσα, παγωτό, λιμάνι.} \\
\texttt{Οι λέξεις της ομάδας του αντιπάλου σας είναι: χαλκός, κύτταρο, σχολείο, πλυντήριο, ζώνη, εργαστήριο, γραμμή, αστυνομία.} \\
\texttt{Οι λέξεις αντιπερισπασμού είναι: αεροπλάνο, λιμουζίνα, προσωπικό, μπότα, κατηγορώ, καρέκλα, μαστίγιο, γεράκι.} \\
\texttt{Οι λέξεις θανάτου είναι: .} \\
            }
        }
    }
    & \\ \\

    \theutterance \stepcounter{utterance}  
    & \multicolumn{4}{p{0.6\linewidth}}{
        \cellcolor[rgb]{0.9,0.9,0.9}{
            \makecell[{{p{\linewidth}}}]{
                \texttt{\tiny{[P1$\rangle$GM]}}
                \texttt{ΥΠΟΔΕΙΞΗ: νήμα} \\
\texttt{ΣΤΟΧΟΙ: χιονάνθρωπος, κάλτσα, συμβολοσειρά, σύνδεσμος} \\
            }
        }
    }
    & & \\ \\

    \theutterance \stepcounter{utterance}  
    & & & \multicolumn{2}{p{0.3\linewidth}}{
        \cellcolor[rgb]{0.9,0.9,0.9}{
            \makecell[{{p{\linewidth}}}]{
                \texttt{\tiny{[GM$|$GM]}}
                \texttt{νήμα} \\
            }
        }
    }
    & & \\ \\

    \theutterance \stepcounter{utterance}  
    & & & \multicolumn{2}{p{0.3\linewidth}}{
        \cellcolor[rgb]{0.9,0.9,0.9}{
            \makecell[{{p{\linewidth}}}]{
                \texttt{\tiny{[GM$|$GM]}}
                ['χιονάνθρωπος', 'κάλτσα', 'συμβολοσειρά', 'σύνδεσμος']
            }
        }
    }
    & & \\ \\

    \theutterance \stepcounter{utterance}  
    & & & \multicolumn{2}{p{0.3\linewidth}}{
        \cellcolor[rgb]{0.9,0.9,0.9}{
            \makecell[{{p{\linewidth}}}]{
                \texttt{\tiny{[GM$|$GM]}}
                {'word': 'χιονάνθρωπος', 'assignment': 'team'}
            }
        }
    }
    & & \\ \\

    \theutterance \stepcounter{utterance}  
    & & & \multicolumn{2}{p{0.3\linewidth}}{
        \cellcolor[rgb]{0.9,0.9,0.9}{
            \makecell[{{p{\linewidth}}}]{
                \texttt{\tiny{[GM$|$GM]}}
                {'word': 'κάλτσα', 'assignment': 'team'}
            }
        }
    }
    & & \\ \\

    \theutterance \stepcounter{utterance}  
    & & & \multicolumn{2}{p{0.3\linewidth}}{
        \cellcolor[rgb]{0.9,0.9,0.9}{
            \makecell[{{p{\linewidth}}}]{
                \texttt{\tiny{[GM$|$GM]}}
                {'word': 'συμβολοσειρά', 'assignment': 'team'}
            }
        }
    }
    & & \\ \\

    \theutterance \stepcounter{utterance}  
    & & & \multicolumn{2}{p{0.3\linewidth}}{
        \cellcolor[rgb]{0.9,0.9,0.9}{
            \makecell[{{p{\linewidth}}}]{
                \texttt{\tiny{[GM$|$GM]}}
                {'word': 'σύνδεσμος', 'assignment': 'team'}
            }
        }
    }
    & & \\ \\

    \theutterance \stepcounter{utterance}  
    & & \multicolumn{4}{p{0.6\linewidth}}{
        \cellcolor[rgb]{0.9,0.9,0.9}{
            \makecell[{{p{\linewidth}}}]{
                \texttt{\tiny{[GM$\rangle$P2]}}
                \texttt{Δώστε μια λίστα χωρισμένη με κόμματα με έως και 4 λέξεις από την παρακάτω λίστα που σχετίζονται καλύτερα ή συνδέονται στενότερα με τη λέξη 'νήμα'. Ξεκινήστε πάντα τη λίστα των εικασιών σας με ΔΟΚΙΜΗ:  και μην συμπεριλάβετε κανένα άλλο κείμενο στην απάντησή σας.} \\
\\ 
\texttt{συμβολοσειρά, μάτι, πλυντήριο, μπότα, χιονάνθρωπος, ζώνη, μαστίγιο, αεροπλάνο, γραμμή, λιμάνι, γεράκι, γκρεμός, προσωπικό, χάρη, σύνδεσμος, εργαστήριο, καρέκλα, αστυνομία, παγωτό, κατηγορώ, κύτταρο, κάλτσα, χαλκός, σχολείο, λιμουζίνα} \\
            }
        }
    }
    & \\ \\

    \theutterance \stepcounter{utterance}  
    & & & \multicolumn{4}{p{0.6\linewidth}}{
        \cellcolor[rgb]{0.9,0.9,0.9}{
            \makecell[{{p{\linewidth}}}]{
                \texttt{\tiny{[GM$\langle$P2]}}
                \texttt{ΔΟΚΙΜΗ: μάτι, γραμμή, κάλτσα, συμβολοσειρά} \\
            }
        }
    }
     \\ \\

    \theutterance \stepcounter{utterance}  
    & & & \multicolumn{2}{p{0.3\linewidth}}{
        \cellcolor[rgb]{0.9,0.9,0.9}{
            \makecell[{{p{\linewidth}}}]{
                \texttt{\tiny{[GM$|$GM]}}
                ['μάτι', 'γραμμή', 'κάλτσα', 'συμβολοσειρά']
            }
        }
    }
    & & \\ \\

    \theutterance \stepcounter{utterance}  
    & & & \multicolumn{2}{p{0.3\linewidth}}{
        \cellcolor[rgb]{0.9,0.9,0.9}{
            \makecell[{{p{\linewidth}}}]{
                \texttt{\tiny{[GM$|$GM]}}
                {'word': 'μάτι', 'assignment': 'team'}
            }
        }
    }
    & & \\ \\

    \theutterance \stepcounter{utterance}  
    & & & \multicolumn{2}{p{0.3\linewidth}}{
        \cellcolor[rgb]{0.9,0.9,0.9}{
            \makecell[{{p{\linewidth}}}]{
                \texttt{\tiny{[GM$|$GM]}}
                {'word': 'γραμμή', 'assignment': 'opponent'}
            }
        }
    }
    & & \\ \\

    \theutterance \stepcounter{utterance}  
    & & & \multicolumn{2}{p{0.3\linewidth}}{
        \cellcolor[rgb]{0.9,0.9,0.9}{
            \makecell[{{p{\linewidth}}}]{
                \texttt{\tiny{[GM$|$GM]}}
                \texttt{γραμμή} \\
            }
        }
    }
    & & \\ \\

    \theutterance \stepcounter{utterance}  
    & & & \multicolumn{2}{p{0.3\linewidth}}{
        \cellcolor[rgb]{0.9,0.9,0.9}{
            \makecell[{{p{\linewidth}}}]{
                \texttt{\tiny{[GM$|$GM]}}
                {'word': 'χαλκός', 'assignment': 'opponent'}
            }
        }
    }
    & & \\ \\

    \theutterance \stepcounter{utterance}  
    & & & \multicolumn{2}{p{0.3\linewidth}}{
        \cellcolor[rgb]{0.9,0.9,0.9}{
            \makecell[{{p{\linewidth}}}]{
                \texttt{\tiny{[GM$|$GM]}}
                {'hidden': {'team': ['χάρη', 'χιονάνθρωπος', 'σύνδεσμος', 'συμβολοσειρά', 'γκρεμός', 'κάλτσα', 'παγωτό', 'λιμάνι'], 'innocent': ['αεροπλάνο', 'λιμουζίνα', 'προσωπικό', 'μπότα', 'κατηγορώ', 'καρέκλα', 'μαστίγιο', 'γεράκι'], 'opponent': ['κύτταρο', 'σχολείο', 'πλυντήριο', 'ζώνη', 'εργαστήριο', 'αστυνομία'], 'assassin': []}, 'revealed': {'team': {'team': ['μάτι'], 'innocent': [], 'opponent': ['γραμμή'], 'assassin': []}, 'opponent': {'team': [], 'innocent': [], 'opponent': ['χαλκός'], 'assassin': []}}}
            }
        }
    }
    & & \\ \\

    \theutterance \stepcounter{utterance}  
    & & \multicolumn{4}{p{0.6\linewidth}}{
        \cellcolor[rgb]{0.9,0.9,0.9}{
            \makecell[{{p{\linewidth}}}]{
                \texttt{\tiny{[P1$\langle$GM]}}
                \texttt{Η λέξη μάτι μαντεύτηκε σωστά. Η λέξη γραμμή μαντεύτηκε, αλλά είναι μια αντίπαλος λέξη. Η σειρά του συμπαίκτη σου έληξε εκεί.} \\
\texttt{Τώρα δώστε ένα άλλο στοιχείο που αφορά κάποιες από τις υπόλοιπες λέξεις της ομάδας σας και έναν κατάλογο των σχετικών λέξεων{-}στόχων. Θυμηθείτε να ξεκινήσετε το στοιχείο σας με 'ΥΠΟΔΕΙΞΗ: ', να βάλετε μια νέα γραμμή και να ξεκινήσετε τον κατάλογο των λέξεων{-}στόχων που χωρίζονται με κόμμα με 'ΣΤΟΧΟΙ: '. Σημείωση: ορισμένες λέξεις έχουν αφαιρεθεί από τις λίστες σε σύγκριση με προηγούμενες αιτήσεις.} \\
\\ 
\texttt{Οι υπόλοιπες λέξεις της ομάδας σας είναι: χάρη, χιονάνθρωπος, σύνδεσμος, συμβολοσειρά, γκρεμός, κάλτσα, παγωτό, λιμάνι.} \\
\texttt{Οι υπόλοιπες λέξεις για τον αντίπαλό σας είναι: χαλκός, κύτταρο, σχολείο, πλυντήριο, ζώνη, εργαστήριο, αστυνομία.} \\
\texttt{Οι υπόλοιπες λέξεις που αποσπούν την προσοχή είναι: αεροπλάνο, λιμουζίνα, προσωπικό, μπότα, κατηγορώ, καρέκλα, μαστίγιο, γεράκι.} \\
\texttt{Οι υπόλοιπες λέξεις που σκοτώνουν είναι: .} \\
            }
        }
    }
    & \\ \\

    \theutterance \stepcounter{utterance}  
    & \multicolumn{4}{p{0.6\linewidth}}{
        \cellcolor[rgb]{0.9,0.9,0.9}{
            \makecell[{{p{\linewidth}}}]{
                \texttt{\tiny{[P1$\rangle$GM]}}
                \texttt{ΥΠΟΔΕΙΞΗ: χειμώνας} \\
\texttt{ΣΤΟΧΟΙ: χιονάνθρωπος, παγωτό, κάλτσα} \\
            }
        }
    }
    & & \\ \\

    \theutterance \stepcounter{utterance}  
    & & & \multicolumn{2}{p{0.3\linewidth}}{
        \cellcolor[rgb]{0.9,0.9,0.9}{
            \makecell[{{p{\linewidth}}}]{
                \texttt{\tiny{[GM$|$GM]}}
                \texttt{χειμώνας} \\
            }
        }
    }
    & & \\ \\

    \theutterance \stepcounter{utterance}  
    & & & \multicolumn{2}{p{0.3\linewidth}}{
        \cellcolor[rgb]{0.9,0.9,0.9}{
            \makecell[{{p{\linewidth}}}]{
                \texttt{\tiny{[GM$|$GM]}}
                ['χιονάνθρωπος', 'παγωτό', 'κάλτσα']
            }
        }
    }
    & & \\ \\

    \theutterance \stepcounter{utterance}  
    & & & \multicolumn{2}{p{0.3\linewidth}}{
        \cellcolor[rgb]{0.9,0.9,0.9}{
            \makecell[{{p{\linewidth}}}]{
                \texttt{\tiny{[GM$|$GM]}}
                {'word': 'χιονάνθρωπος', 'assignment': 'team'}
            }
        }
    }
    & & \\ \\

    \theutterance \stepcounter{utterance}  
    & & & \multicolumn{2}{p{0.3\linewidth}}{
        \cellcolor[rgb]{0.9,0.9,0.9}{
            \makecell[{{p{\linewidth}}}]{
                \texttt{\tiny{[GM$|$GM]}}
                {'word': 'παγωτό', 'assignment': 'team'}
            }
        }
    }
    & & \\ \\

    \theutterance \stepcounter{utterance}  
    & & & \multicolumn{2}{p{0.3\linewidth}}{
        \cellcolor[rgb]{0.9,0.9,0.9}{
            \makecell[{{p{\linewidth}}}]{
                \texttt{\tiny{[GM$|$GM]}}
                {'word': 'κάλτσα', 'assignment': 'team'}
            }
        }
    }
    & & \\ \\

    \theutterance \stepcounter{utterance}  
    & & \multicolumn{4}{p{0.6\linewidth}}{
        \cellcolor[rgb]{0.9,0.9,0.9}{
            \makecell[{{p{\linewidth}}}]{
                \texttt{\tiny{[GM$\rangle$P2]}}
                \texttt{Η λέξη μάτι μαντεύτηκε σωστά. Η λέξη γραμμή μαντεύτηκε, αλλά είναι μια αντίπαλος λέξη. Η σειρά σου έληξε εκεί.} \\
\texttt{Τώρα δώστε έναν άλλο κατάλογο διαχωρισμένο με κόμματα από τουλάχιστον 1 και έως 3 λέξεις από την ακόλουθη λίστα λέξεων που σχετίζονται καλύτερα ή συνδέονται πιο στενά με τη λέξη 'χειμώνας'. Θυμηθείτε να ξεκινήσετε την απάντησή σας με ΔΟΚΙΜΗ: . Σημείωση: ορισμένες λέξεις έχουν αφαιρεθεί από τη λίστα σε σύγκριση με προηγούμενα αιτήματα.} \\
\\ 
\texttt{συμβολοσειρά, πλυντήριο, μπότα, χιονάνθρωπος, ζώνη, μαστίγιο, αεροπλάνο, λιμάνι, γεράκι, γκρεμός, προσωπικό, χάρη, σύνδεσμος, εργαστήριο, καρέκλα, αστυνομία, παγωτό, κατηγορώ, κύτταρο, κάλτσα, σχολείο, λιμουζίνα} \\
            }
        }
    }
    & \\ \\

    \theutterance \stepcounter{utterance}  
    & & & \multicolumn{4}{p{0.6\linewidth}}{
        \cellcolor[rgb]{0.9,0.9,0.9}{
            \makecell[{{p{\linewidth}}}]{
                \texttt{\tiny{[GM$\langle$P2]}}
                \texttt{ΔΟΚΙΜΗ: χιονάνθρωπος, μπότα, παγωτό} \\
            }
        }
    }
     \\ \\

    \theutterance \stepcounter{utterance}  
    & & & \multicolumn{2}{p{0.3\linewidth}}{
        \cellcolor[rgb]{0.9,0.9,0.9}{
            \makecell[{{p{\linewidth}}}]{
                \texttt{\tiny{[GM$|$GM]}}
                ['χιονάνθρωπος', 'μπότα', 'παγωτό']
            }
        }
    }
    & & \\ \\

    \theutterance \stepcounter{utterance}  
    & & & \multicolumn{2}{p{0.3\linewidth}}{
        \cellcolor[rgb]{0.9,0.9,0.9}{
            \makecell[{{p{\linewidth}}}]{
                \texttt{\tiny{[GM$|$GM]}}
                {'word': 'χιονάνθρωπος', 'assignment': 'team'}
            }
        }
    }
    & & \\ \\

    \theutterance \stepcounter{utterance}  
    & & & \multicolumn{2}{p{0.3\linewidth}}{
        \cellcolor[rgb]{0.9,0.9,0.9}{
            \makecell[{{p{\linewidth}}}]{
                \texttt{\tiny{[GM$|$GM]}}
                {'word': 'χιονάνθρωπος', 'assignment': 'team'}
            }
        }
    }
    & & \\ \\

    \theutterance \stepcounter{utterance}  
    & & & \multicolumn{2}{p{0.3\linewidth}}{
        \cellcolor[rgb]{0.9,0.9,0.9}{
            \makecell[{{p{\linewidth}}}]{
                \texttt{\tiny{[GM$|$GM]}}
                {'word': 'μπότα', 'assignment': 'innocent'}
            }
        }
    }
    & & \\ \\

    \theutterance \stepcounter{utterance}  
    & & & \multicolumn{2}{p{0.3\linewidth}}{
        \cellcolor[rgb]{0.9,0.9,0.9}{
            \makecell[{{p{\linewidth}}}]{
                \texttt{\tiny{[GM$|$GM]}}
                \texttt{μπότα} \\
            }
        }
    }
    & & \\ \\

    \theutterance \stepcounter{utterance}  
    & & & \multicolumn{2}{p{0.3\linewidth}}{
        \cellcolor[rgb]{0.9,0.9,0.9}{
            \makecell[{{p{\linewidth}}}]{
                \texttt{\tiny{[GM$|$GM]}}
                {'word': 'κύτταρο', 'assignment': 'opponent'}
            }
        }
    }
    & & \\ \\

    \theutterance \stepcounter{utterance}  
    & & & \multicolumn{2}{p{0.3\linewidth}}{
        \cellcolor[rgb]{0.9,0.9,0.9}{
            \makecell[{{p{\linewidth}}}]{
                \texttt{\tiny{[GM$|$GM]}}
                {'hidden': {'team': ['χάρη', 'σύνδεσμος', 'συμβολοσειρά', 'γκρεμός', 'κάλτσα', 'παγωτό', 'λιμάνι'], 'innocent': ['αεροπλάνο', 'λιμουζίνα', 'προσωπικό', 'κατηγορώ', 'καρέκλα', 'μαστίγιο', 'γεράκι'], 'opponent': ['σχολείο', 'πλυντήριο', 'ζώνη', 'εργαστήριο', 'αστυνομία'], 'assassin': []}, 'revealed': {'team': {'team': ['μάτι', 'χιονάνθρωπος'], 'innocent': ['μπότα'], 'opponent': ['γραμμή'], 'assassin': []}, 'opponent': {'team': [], 'innocent': [], 'opponent': ['χαλκός', 'κύτταρο'], 'assassin': []}}}
            }
        }
    }
    & & \\ \\

    \theutterance \stepcounter{utterance}  
    & & \multicolumn{4}{p{0.6\linewidth}}{
        \cellcolor[rgb]{0.9,0.9,0.9}{
            \makecell[{{p{\linewidth}}}]{
                \texttt{\tiny{[P1$\langle$GM]}}
                \texttt{Η λέξη χιονάνθρωπος μαντεύτηκε σωστά. Η λέξη μπότα μαντεύτηκε, αλλά είναι μια αθώος λέξη. Η σειρά του συμπαίκτη σου έληξε εκεί.} \\
\texttt{Τώρα δώστε ένα άλλο στοιχείο που αφορά κάποιες από τις υπόλοιπες λέξεις της ομάδας σας και έναν κατάλογο των σχετικών λέξεων{-}στόχων. Θυμηθείτε να ξεκινήσετε το στοιχείο σας με 'ΥΠΟΔΕΙΞΗ: ', να βάλετε μια νέα γραμμή και να ξεκινήσετε τον κατάλογο των λέξεων{-}στόχων που χωρίζονται με κόμμα με 'ΣΤΟΧΟΙ: '. Σημείωση: ορισμένες λέξεις έχουν αφαιρεθεί από τις λίστες σε σύγκριση με προηγούμενες αιτήσεις.} \\
\\ 
\texttt{Οι υπόλοιπες λέξεις της ομάδας σας είναι: χάρη, σύνδεσμος, συμβολοσειρά, γκρεμός, κάλτσα, παγωτό, λιμάνι.} \\
\texttt{Οι υπόλοιπες λέξεις για τον αντίπαλό σας είναι: κύτταρο, σχολείο, πλυντήριο, ζώνη, εργαστήριο, αστυνομία.} \\
\texttt{Οι υπόλοιπες λέξεις που αποσπούν την προσοχή είναι: αεροπλάνο, λιμουζίνα, προσωπικό, κατηγορώ, καρέκλα, μαστίγιο, γεράκι.} \\
\texttt{Οι υπόλοιπες λέξεις που σκοτώνουν είναι: .} \\
            }
        }
    }
    & \\ \\

    \theutterance \stepcounter{utterance}  
    & \multicolumn{4}{p{0.6\linewidth}}{
        \cellcolor[rgb]{0.9,0.9,0.9}{
            \makecell[{{p{\linewidth}}}]{
                \texttt{\tiny{[P1$\rangle$GM]}}
                \texttt{ΥΠΟΔΕΙΞΗ: πληροφορική} \\
\texttt{ΣΤΟΧΟΙ: σύνδεσμος, συμβολοσειρά} \\
            }
        }
    }
    & & \\ \\

    \theutterance \stepcounter{utterance}  
    & & & \multicolumn{2}{p{0.3\linewidth}}{
        \cellcolor[rgb]{0.9,0.9,0.9}{
            \makecell[{{p{\linewidth}}}]{
                \texttt{\tiny{[GM$|$GM]}}
                \texttt{πληροφορική} \\
            }
        }
    }
    & & \\ \\

    \theutterance \stepcounter{utterance}  
    & & & \multicolumn{2}{p{0.3\linewidth}}{
        \cellcolor[rgb]{0.9,0.9,0.9}{
            \makecell[{{p{\linewidth}}}]{
                \texttt{\tiny{[GM$|$GM]}}
                ['σύνδεσμος', 'συμβολοσειρά']
            }
        }
    }
    & & \\ \\

    \theutterance \stepcounter{utterance}  
    & & & \multicolumn{2}{p{0.3\linewidth}}{
        \cellcolor[rgb]{0.9,0.9,0.9}{
            \makecell[{{p{\linewidth}}}]{
                \texttt{\tiny{[GM$|$GM]}}
                {'word': 'σύνδεσμος', 'assignment': 'team'}
            }
        }
    }
    & & \\ \\

    \theutterance \stepcounter{utterance}  
    & & & \multicolumn{2}{p{0.3\linewidth}}{
        \cellcolor[rgb]{0.9,0.9,0.9}{
            \makecell[{{p{\linewidth}}}]{
                \texttt{\tiny{[GM$|$GM]}}
                {'word': 'συμβολοσειρά', 'assignment': 'team'}
            }
        }
    }
    & & \\ \\

    \theutterance \stepcounter{utterance}  
    & & \multicolumn{4}{p{0.6\linewidth}}{
        \cellcolor[rgb]{0.9,0.9,0.9}{
            \makecell[{{p{\linewidth}}}]{
                \texttt{\tiny{[GM$\rangle$P2]}}
                \texttt{Η λέξη χιονάνθρωπος μαντεύτηκε σωστά. Η λέξη μπότα μαντεύτηκε, αλλά είναι μια αθώος λέξη. Η σειρά σου έληξε εκεί.} \\
\texttt{Τώρα δώστε έναν άλλο κατάλογο διαχωρισμένο με κόμματα από τουλάχιστον 1 και έως 2 λέξεις από την ακόλουθη λίστα λέξεων που σχετίζονται καλύτερα ή συνδέονται πιο στενά με τη λέξη 'πληροφορική'. Θυμηθείτε να ξεκινήσετε την απάντησή σας με ΔΟΚΙΜΗ: . Σημείωση: ορισμένες λέξεις έχουν αφαιρεθεί από τη λίστα σε σύγκριση με προηγούμενα αιτήματα.} \\
\\ 
\texttt{συμβολοσειρά, πλυντήριο, ζώνη, μαστίγιο, αεροπλάνο, λιμάνι, γεράκι, γκρεμός, προσωπικό, χάρη, σύνδεσμος, εργαστήριο, καρέκλα, αστυνομία, παγωτό, κατηγορώ, κάλτσα, σχολείο, λιμουζίνα} \\
            }
        }
    }
    & \\ \\

    \theutterance \stepcounter{utterance}  
    & & & \multicolumn{4}{p{0.6\linewidth}}{
        \cellcolor[rgb]{0.9,0.9,0.9}{
            \makecell[{{p{\linewidth}}}]{
                \texttt{\tiny{[GM$\langle$P2]}}
                \texttt{ΔΟΚΙΜΗ: συμβολοσειρά, σύνδεσμος} \\
            }
        }
    }
     \\ \\

    \theutterance \stepcounter{utterance}  
    & & & \multicolumn{2}{p{0.3\linewidth}}{
        \cellcolor[rgb]{0.9,0.9,0.9}{
            \makecell[{{p{\linewidth}}}]{
                \texttt{\tiny{[GM$|$GM]}}
                ['συμβολοσειρά', 'σύνδεσμος']
            }
        }
    }
    & & \\ \\

    \theutterance \stepcounter{utterance}  
    & & & \multicolumn{2}{p{0.3\linewidth}}{
        \cellcolor[rgb]{0.9,0.9,0.9}{
            \makecell[{{p{\linewidth}}}]{
                \texttt{\tiny{[GM$|$GM]}}
                {'word': 'συμβολοσειρά', 'assignment': 'team'}
            }
        }
    }
    & & \\ \\

    \theutterance \stepcounter{utterance}  
    & & & \multicolumn{2}{p{0.3\linewidth}}{
        \cellcolor[rgb]{0.9,0.9,0.9}{
            \makecell[{{p{\linewidth}}}]{
                \texttt{\tiny{[GM$|$GM]}}
                {'word': 'συμβολοσειρά', 'assignment': 'team'}
            }
        }
    }
    & & \\ \\

    \theutterance \stepcounter{utterance}  
    & & & \multicolumn{2}{p{0.3\linewidth}}{
        \cellcolor[rgb]{0.9,0.9,0.9}{
            \makecell[{{p{\linewidth}}}]{
                \texttt{\tiny{[GM$|$GM]}}
                {'word': 'σύνδεσμος', 'assignment': 'team'}
            }
        }
    }
    & & \\ \\

    \theutterance \stepcounter{utterance}  
    & & & \multicolumn{2}{p{0.3\linewidth}}{
        \cellcolor[rgb]{0.9,0.9,0.9}{
            \makecell[{{p{\linewidth}}}]{
                \texttt{\tiny{[GM$|$GM]}}
                {'word': 'σύνδεσμος', 'assignment': 'team'}
            }
        }
    }
    & & \\ \\

    \theutterance \stepcounter{utterance}  
    & & & \multicolumn{2}{p{0.3\linewidth}}{
        \cellcolor[rgb]{0.9,0.9,0.9}{
            \makecell[{{p{\linewidth}}}]{
                \texttt{\tiny{[GM$|$GM]}}
                {'word': 'πλυντήριο', 'assignment': 'opponent'}
            }
        }
    }
    & & \\ \\

    \theutterance \stepcounter{utterance}  
    & & & \multicolumn{2}{p{0.3\linewidth}}{
        \cellcolor[rgb]{0.9,0.9,0.9}{
            \makecell[{{p{\linewidth}}}]{
                \texttt{\tiny{[GM$|$GM]}}
                {'hidden': {'team': ['χάρη', 'γκρεμός', 'κάλτσα', 'παγωτό', 'λιμάνι'], 'innocent': ['αεροπλάνο', 'λιμουζίνα', 'προσωπικό', 'κατηγορώ', 'καρέκλα', 'μαστίγιο', 'γεράκι'], 'opponent': ['σχολείο', 'ζώνη', 'εργαστήριο', 'αστυνομία'], 'assassin': []}, 'revealed': {'team': {'team': ['μάτι', 'χιονάνθρωπος', 'συμβολοσειρά', 'σύνδεσμος'], 'innocent': ['μπότα'], 'opponent': ['γραμμή'], 'assassin': []}, 'opponent': {'team': [], 'innocent': [], 'opponent': ['χαλκός', 'κύτταρο', 'πλυντήριο'], 'assassin': []}}}
            }
        }
    }
    & & \\ \\

    \theutterance \stepcounter{utterance}  
    & & \multicolumn{4}{p{0.6\linewidth}}{
        \cellcolor[rgb]{0.9,0.9,0.9}{
            \makecell[{{p{\linewidth}}}]{
                \texttt{\tiny{[P1$\langle$GM]}}
                \texttt{Οι λέξεις συμβολοσειρά, σύνδεσμος μαντεύτηκαν σωστά. Η σειρά του συμπαίκτη σου έληξε εκεί.} \\
\texttt{Τώρα δώστε ένα άλλο στοιχείο που αφορά κάποιες από τις υπόλοιπες λέξεις της ομάδας σας και έναν κατάλογο των σχετικών λέξεων{-}στόχων. Θυμηθείτε να ξεκινήσετε το στοιχείο σας με 'ΥΠΟΔΕΙΞΗ: ', να βάλετε μια νέα γραμμή και να ξεκινήσετε τον κατάλογο των λέξεων{-}στόχων που χωρίζονται με κόμμα με 'ΣΤΟΧΟΙ: '. Σημείωση: ορισμένες λέξεις έχουν αφαιρεθεί από τις λίστες σε σύγκριση με προηγούμενες αιτήσεις.} \\
\\ 
\texttt{Οι υπόλοιπες λέξεις της ομάδας σας είναι: χάρη, γκρεμός, κάλτσα, παγωτό, λιμάνι.} \\
\texttt{Οι υπόλοιπες λέξεις για τον αντίπαλό σας είναι: σχολείο, πλυντήριο, ζώνη, εργαστήριο, αστυνομία.} \\
\texttt{Οι υπόλοιπες λέξεις που αποσπούν την προσοχή είναι: αεροπλάνο, λιμουζίνα, προσωπικό, κατηγορώ, καρέκλα, μαστίγιο, γεράκι.} \\
\texttt{Οι υπόλοιπες λέξεις που σκοτώνουν είναι: .} \\
            }
        }
    }
    & \\ \\

    \theutterance \stepcounter{utterance}  
    & \multicolumn{4}{p{0.6\linewidth}}{
        \cellcolor[rgb]{0.9,0.9,0.9}{
            \makecell[{{p{\linewidth}}}]{
                \texttt{\tiny{[P1$\rangle$GM]}}
                \texttt{ΥΠΟΔΕΙΞΗ: παραλιακό} \\
\texttt{ΣΤΟΧΟΙ: γκρεμός, λιμάνι, παγωτό} \\
            }
        }
    }
    & & \\ \\

    \theutterance \stepcounter{utterance}  
    & & & \multicolumn{2}{p{0.3\linewidth}}{
        \cellcolor[rgb]{0.9,0.9,0.9}{
            \makecell[{{p{\linewidth}}}]{
                \texttt{\tiny{[GM$|$GM]}}
                \texttt{παραλιακό} \\
            }
        }
    }
    & & \\ \\

    \theutterance \stepcounter{utterance}  
    & & & \multicolumn{2}{p{0.3\linewidth}}{
        \cellcolor[rgb]{0.9,0.9,0.9}{
            \makecell[{{p{\linewidth}}}]{
                \texttt{\tiny{[GM$|$GM]}}
                ['γκρεμός', 'λιμάνι', 'παγωτό']
            }
        }
    }
    & & \\ \\

    \theutterance \stepcounter{utterance}  
    & & & \multicolumn{2}{p{0.3\linewidth}}{
        \cellcolor[rgb]{0.9,0.9,0.9}{
            \makecell[{{p{\linewidth}}}]{
                \texttt{\tiny{[GM$|$GM]}}
                {'word': 'γκρεμός', 'assignment': 'team'}
            }
        }
    }
    & & \\ \\

    \theutterance \stepcounter{utterance}  
    & & & \multicolumn{2}{p{0.3\linewidth}}{
        \cellcolor[rgb]{0.9,0.9,0.9}{
            \makecell[{{p{\linewidth}}}]{
                \texttt{\tiny{[GM$|$GM]}}
                {'word': 'λιμάνι', 'assignment': 'team'}
            }
        }
    }
    & & \\ \\

    \theutterance \stepcounter{utterance}  
    & & & \multicolumn{2}{p{0.3\linewidth}}{
        \cellcolor[rgb]{0.9,0.9,0.9}{
            \makecell[{{p{\linewidth}}}]{
                \texttt{\tiny{[GM$|$GM]}}
                {'word': 'παγωτό', 'assignment': 'team'}
            }
        }
    }
    & & \\ \\

    \theutterance \stepcounter{utterance}  
    & & \multicolumn{4}{p{0.6\linewidth}}{
        \cellcolor[rgb]{0.9,0.9,0.9}{
            \makecell[{{p{\linewidth}}}]{
                \texttt{\tiny{[GM$\rangle$P2]}}
                \texttt{Οι λέξεις συμβολοσειρά, σύνδεσμος μαντεύτηκαν σωστά. Η σειρά σου έληξε εκεί.} \\
\texttt{Τώρα δώστε έναν άλλο κατάλογο διαχωρισμένο με κόμματα από τουλάχιστον 1 και έως 3 λέξεις από την ακόλουθη λίστα λέξεων που σχετίζονται καλύτερα ή συνδέονται πιο στενά με τη λέξη 'παραλιακό'. Θυμηθείτε να ξεκινήσετε την απάντησή σας με ΔΟΚΙΜΗ: . Σημείωση: ορισμένες λέξεις έχουν αφαιρεθεί από τη λίστα σε σύγκριση με προηγούμενα αιτήματα.} \\
\\ 
\texttt{ζώνη, μαστίγιο, αεροπλάνο, λιμάνι, γεράκι, γκρεμός, προσωπικό, χάρη, εργαστήριο, καρέκλα, αστυνομία, παγωτό, κατηγορώ, κάλτσα, σχολείο, λιμουζίνα} \\
            }
        }
    }
    & \\ \\

    \theutterance \stepcounter{utterance}  
    & & & \multicolumn{4}{p{0.6\linewidth}}{
        \cellcolor[rgb]{0.9,0.9,0.9}{
            \makecell[{{p{\linewidth}}}]{
                \texttt{\tiny{[GM$\langle$P2]}}
                \texttt{ΔΟΚΙΜΗ: λιμάνι, γκρεμός, αεροπλάνο} \\
            }
        }
    }
     \\ \\

    \theutterance \stepcounter{utterance}  
    & & & \multicolumn{2}{p{0.3\linewidth}}{
        \cellcolor[rgb]{0.9,0.9,0.9}{
            \makecell[{{p{\linewidth}}}]{
                \texttt{\tiny{[GM$|$GM]}}
                ['λιμάνι', 'γκρεμός', 'αεροπλάνο']
            }
        }
    }
    & & \\ \\

    \theutterance \stepcounter{utterance}  
    & & & \multicolumn{2}{p{0.3\linewidth}}{
        \cellcolor[rgb]{0.9,0.9,0.9}{
            \makecell[{{p{\linewidth}}}]{
                \texttt{\tiny{[GM$|$GM]}}
                {'word': 'λιμάνι', 'assignment': 'team'}
            }
        }
    }
    & & \\ \\

    \theutterance \stepcounter{utterance}  
    & & & \multicolumn{2}{p{0.3\linewidth}}{
        \cellcolor[rgb]{0.9,0.9,0.9}{
            \makecell[{{p{\linewidth}}}]{
                \texttt{\tiny{[GM$|$GM]}}
                {'word': 'λιμάνι', 'assignment': 'team'}
            }
        }
    }
    & & \\ \\

    \theutterance \stepcounter{utterance}  
    & & & \multicolumn{2}{p{0.3\linewidth}}{
        \cellcolor[rgb]{0.9,0.9,0.9}{
            \makecell[{{p{\linewidth}}}]{
                \texttt{\tiny{[GM$|$GM]}}
                {'word': 'γκρεμός', 'assignment': 'team'}
            }
        }
    }
    & & \\ \\

    \theutterance \stepcounter{utterance}  
    & & & \multicolumn{2}{p{0.3\linewidth}}{
        \cellcolor[rgb]{0.9,0.9,0.9}{
            \makecell[{{p{\linewidth}}}]{
                \texttt{\tiny{[GM$|$GM]}}
                {'word': 'γκρεμός', 'assignment': 'team'}
            }
        }
    }
    & & \\ \\

    \theutterance \stepcounter{utterance}  
    & & & \multicolumn{2}{p{0.3\linewidth}}{
        \cellcolor[rgb]{0.9,0.9,0.9}{
            \makecell[{{p{\linewidth}}}]{
                \texttt{\tiny{[GM$|$GM]}}
                {'word': 'αεροπλάνο', 'assignment': 'innocent'}
            }
        }
    }
    & & \\ \\

    \theutterance \stepcounter{utterance}  
    & & & \multicolumn{2}{p{0.3\linewidth}}{
        \cellcolor[rgb]{0.9,0.9,0.9}{
            \makecell[{{p{\linewidth}}}]{
                \texttt{\tiny{[GM$|$GM]}}
                \texttt{αεροπλάνο} \\
            }
        }
    }
    & & \\ \\

    \theutterance \stepcounter{utterance}  
    & & & \multicolumn{2}{p{0.3\linewidth}}{
        \cellcolor[rgb]{0.9,0.9,0.9}{
            \makecell[{{p{\linewidth}}}]{
                \texttt{\tiny{[GM$|$GM]}}
                {'word': 'ζώνη', 'assignment': 'opponent'}
            }
        }
    }
    & & \\ \\

    \theutterance \stepcounter{utterance}  
    & & & \multicolumn{2}{p{0.3\linewidth}}{
        \cellcolor[rgb]{0.9,0.9,0.9}{
            \makecell[{{p{\linewidth}}}]{
                \texttt{\tiny{[GM$|$GM]}}
                {'hidden': {'team': ['χάρη', 'κάλτσα', 'παγωτό'], 'innocent': ['λιμουζίνα', 'προσωπικό', 'κατηγορώ', 'καρέκλα', 'μαστίγιο', 'γεράκι'], 'opponent': ['σχολείο', 'εργαστήριο', 'αστυνομία'], 'assassin': []}, 'revealed': {'team': {'team': ['μάτι', 'χιονάνθρωπος', 'συμβολοσειρά', 'σύνδεσμος', 'λιμάνι', 'γκρεμός'], 'innocent': ['μπότα', 'αεροπλάνο'], 'opponent': ['γραμμή'], 'assassin': []}, 'opponent': {'team': [], 'innocent': [], 'opponent': ['χαλκός', 'κύτταρο', 'πλυντήριο', 'ζώνη'], 'assassin': []}}}
            }
        }
    }
    & & \\ \\

    \theutterance \stepcounter{utterance}  
    & & \multicolumn{4}{p{0.6\linewidth}}{
        \cellcolor[rgb]{0.9,0.9,0.9}{
            \makecell[{{p{\linewidth}}}]{
                \texttt{\tiny{[P1$\langle$GM]}}
                \texttt{Οι λέξεις λιμάνι, γκρεμός μαντεύτηκαν σωστά. Η λέξη αεροπλάνο μαντεύτηκε, αλλά είναι μια αθώος λέξη. Η σειρά του συμπαίκτη σου έληξε εκεί.} \\
\texttt{Τώρα δώστε ένα άλλο στοιχείο που αφορά κάποιες από τις υπόλοιπες λέξεις της ομάδας σας και έναν κατάλογο των σχετικών λέξεων{-}στόχων. Θυμηθείτε να ξεκινήσετε το στοιχείο σας με 'ΥΠΟΔΕΙΞΗ: ', να βάλετε μια νέα γραμμή και να ξεκινήσετε τον κατάλογο των λέξεων{-}στόχων που χωρίζονται με κόμμα με 'ΣΤΟΧΟΙ: '. Σημείωση: ορισμένες λέξεις έχουν αφαιρεθεί από τις λίστες σε σύγκριση με προηγούμενες αιτήσεις.} \\
\\ 
\texttt{Οι υπόλοιπες λέξεις της ομάδας σας είναι: χάρη, κάλτσα, παγωτό.} \\
\texttt{Οι υπόλοιπες λέξεις για τον αντίπαλό σας είναι: σχολείο, ζώνη, εργαστήριο, αστυνομία.} \\
\texttt{Οι υπόλοιπες λέξεις που αποσπούν την προσοχή είναι: λιμουζίνα, προσωπικό, κατηγορώ, καρέκλα, μαστίγιο, γεράκι.} \\
\texttt{Οι υπόλοιπες λέξεις που σκοτώνουν είναι: .} \\
            }
        }
    }
    & \\ \\

    \theutterance \stepcounter{utterance}  
    & \multicolumn{4}{p{0.6\linewidth}}{
        \cellcolor[rgb]{0.9,0.9,0.9}{
            \makecell[{{p{\linewidth}}}]{
                \texttt{\tiny{[P1$\rangle$GM]}}
                \texttt{ΥΠΟΔΕΙΞΗ: δώρο} \\
\texttt{ΣΤΟΧΟΙ: χάρη, κάλτσα, παγωτό} \\
            }
        }
    }
    & & \\ \\

    \theutterance \stepcounter{utterance}  
    & & & \multicolumn{2}{p{0.3\linewidth}}{
        \cellcolor[rgb]{0.9,0.9,0.9}{
            \makecell[{{p{\linewidth}}}]{
                \texttt{\tiny{[GM$|$GM]}}
                \texttt{δώρο} \\
            }
        }
    }
    & & \\ \\

    \theutterance \stepcounter{utterance}  
    & & & \multicolumn{2}{p{0.3\linewidth}}{
        \cellcolor[rgb]{0.9,0.9,0.9}{
            \makecell[{{p{\linewidth}}}]{
                \texttt{\tiny{[GM$|$GM]}}
                ['χάρη', 'κάλτσα', 'παγωτό']
            }
        }
    }
    & & \\ \\

    \theutterance \stepcounter{utterance}  
    & & & \multicolumn{2}{p{0.3\linewidth}}{
        \cellcolor[rgb]{0.9,0.9,0.9}{
            \makecell[{{p{\linewidth}}}]{
                \texttt{\tiny{[GM$|$GM]}}
                {'word': 'χάρη', 'assignment': 'team'}
            }
        }
    }
    & & \\ \\

    \theutterance \stepcounter{utterance}  
    & & & \multicolumn{2}{p{0.3\linewidth}}{
        \cellcolor[rgb]{0.9,0.9,0.9}{
            \makecell[{{p{\linewidth}}}]{
                \texttt{\tiny{[GM$|$GM]}}
                {'word': 'κάλτσα', 'assignment': 'team'}
            }
        }
    }
    & & \\ \\

    \theutterance \stepcounter{utterance}  
    & & & \multicolumn{2}{p{0.3\linewidth}}{
        \cellcolor[rgb]{0.9,0.9,0.9}{
            \makecell[{{p{\linewidth}}}]{
                \texttt{\tiny{[GM$|$GM]}}
                {'word': 'παγωτό', 'assignment': 'team'}
            }
        }
    }
    & & \\ \\

    \theutterance \stepcounter{utterance}  
    & & \multicolumn{4}{p{0.6\linewidth}}{
        \cellcolor[rgb]{0.9,0.9,0.9}{
            \makecell[{{p{\linewidth}}}]{
                \texttt{\tiny{[GM$\rangle$P2]}}
                \texttt{Οι λέξεις λιμάνι, γκρεμός μαντεύτηκαν σωστά. Η λέξη αεροπλάνο μαντεύτηκε, αλλά είναι μια αθώος λέξη. Η σειρά σου έληξε εκεί.} \\
\texttt{Τώρα δώστε έναν άλλο κατάλογο διαχωρισμένο με κόμματα από τουλάχιστον 1 και έως 3 λέξεις από την ακόλουθη λίστα λέξεων που σχετίζονται καλύτερα ή συνδέονται πιο στενά με τη λέξη 'δώρο'. Θυμηθείτε να ξεκινήσετε την απάντησή σας με ΔΟΚΙΜΗ: . Σημείωση: ορισμένες λέξεις έχουν αφαιρεθεί από τη λίστα σε σύγκριση με προηγούμενα αιτήματα.} \\
\\ 
\texttt{μαστίγιο, γεράκι, προσωπικό, χάρη, εργαστήριο, καρέκλα, αστυνομία, παγωτό, κατηγορώ, κάλτσα, σχολείο, λιμουζίνα} \\
            }
        }
    }
    & \\ \\

    \theutterance \stepcounter{utterance}  
    & & & \multicolumn{4}{p{0.6\linewidth}}{
        \cellcolor[rgb]{0.9,0.9,0.9}{
            \makecell[{{p{\linewidth}}}]{
                \texttt{\tiny{[GM$\langle$P2]}}
                \texttt{ΔΟΚΙΜΗ: χάρη, παγωτό, κάλτσα} \\
            }
        }
    }
     \\ \\

    \theutterance \stepcounter{utterance}  
    & & & \multicolumn{2}{p{0.3\linewidth}}{
        \cellcolor[rgb]{0.9,0.9,0.9}{
            \makecell[{{p{\linewidth}}}]{
                \texttt{\tiny{[GM$|$GM]}}
                ['χάρη', 'παγωτό', 'κάλτσα']
            }
        }
    }
    & & \\ \\

    \theutterance \stepcounter{utterance}  
    & & & \multicolumn{2}{p{0.3\linewidth}}{
        \cellcolor[rgb]{0.9,0.9,0.9}{
            \makecell[{{p{\linewidth}}}]{
                \texttt{\tiny{[GM$|$GM]}}
                {'word': 'χάρη', 'assignment': 'team'}
            }
        }
    }
    & & \\ \\

    \theutterance \stepcounter{utterance}  
    & & & \multicolumn{2}{p{0.3\linewidth}}{
        \cellcolor[rgb]{0.9,0.9,0.9}{
            \makecell[{{p{\linewidth}}}]{
                \texttt{\tiny{[GM$|$GM]}}
                {'word': 'χάρη', 'assignment': 'team'}
            }
        }
    }
    & & \\ \\

    \theutterance \stepcounter{utterance}  
    & & & \multicolumn{2}{p{0.3\linewidth}}{
        \cellcolor[rgb]{0.9,0.9,0.9}{
            \makecell[{{p{\linewidth}}}]{
                \texttt{\tiny{[GM$|$GM]}}
                {'word': 'παγωτό', 'assignment': 'team'}
            }
        }
    }
    & & \\ \\

    \theutterance \stepcounter{utterance}  
    & & & \multicolumn{2}{p{0.3\linewidth}}{
        \cellcolor[rgb]{0.9,0.9,0.9}{
            \makecell[{{p{\linewidth}}}]{
                \texttt{\tiny{[GM$|$GM]}}
                {'word': 'παγωτό', 'assignment': 'team'}
            }
        }
    }
    & & \\ \\

    \theutterance \stepcounter{utterance}  
    & & & \multicolumn{2}{p{0.3\linewidth}}{
        \cellcolor[rgb]{0.9,0.9,0.9}{
            \makecell[{{p{\linewidth}}}]{
                \texttt{\tiny{[GM$|$GM]}}
                {'word': 'κάλτσα', 'assignment': 'team'}
            }
        }
    }
    & & \\ \\

    \theutterance \stepcounter{utterance}  
    & & & \multicolumn{2}{p{0.3\linewidth}}{
        \cellcolor[rgb]{0.9,0.9,0.9}{
            \makecell[{{p{\linewidth}}}]{
                \texttt{\tiny{[GM$|$GM]}}
                {'word': 'κάλτσα', 'assignment': 'team'}
            }
        }
    }
    & & \\ \\

    \theutterance \stepcounter{utterance}  
    & & & \multicolumn{2}{p{0.3\linewidth}}{
        \cellcolor[rgb]{0.9,0.9,0.9}{
            \makecell[{{p{\linewidth}}}]{
                \texttt{\tiny{[GM$|$GM]}}
                \texttt{team has won} \\
            }
        }
    }
    & & \\ \\

\end{supertabular}
}

\end{document}
